\documentclass[12pt]{article}
\usepackage[utf8]{inputenc}
\usepackage[T1]{fontenc}
\usepackage{amsmath,amssymb,amscd}
\usepackage{geometry}
\geometry{margin=2.5cm}

\begin{document}

\section*{Curs 3 -- Pagina 4}

\noindent Fie $(\mathbb{Z},+)$ = gr.\ comut., $(M,\cdot,e)$ = monoid.

\noindent Avem $f : \mathbb{Z} \to M$, $f(k) = a^k$

\[
f(h+k) = f(h)\cdot f(k),
\]
adică $a^{h+k} = a^h \cdot a^k$.

\[
\begin{aligned}
h&=m, & k&=n\\
h&=-m, & k&=-n
\end{aligned}
\qquad m,n\in\mathbb{N}^*.
\]

\[
a^{-m}\cdot a^{-n} = (a^{-1})^{m}\cdot(a^{-1})^{n} = (a^{-1})^{m+n} = a^{h+k}
\]
\noindent\textit{(ultimul termen e neclar în original -- ar putea fi $a^{h+k}$ sau $a^h\cdot a^k$; te rog verifică.)}

\bigskip
\noindent\underline{\textbf{Lege de compoziție indusă pe o parte stabilă}} (respectiv pe o mulțime finită)

\bigskip
\noindent\textbf{Def.} Fie $(S,\varphi)$; $\varphi : S\times S \to S$, $S \neq \varnothing$.

\noindent O submulțime $T \subseteq S$ s.n. \textbf{parte stabilă} a lui $S$ în raport cu $\varphi$ dacă $(\forall)\, x,y \in T \Rightarrow \varphi(x,y) \in T$.

\[
\begin{CD}
S\times S @>\varphi>> S\\
@AAA @AAA\\
T\times T @>\varphi'>> T
\end{CD}
\]
\noindent $(x,y)\in T \Rightarrow \varphi'(x,y) \overset{\text{def}}{=} \varphi(x,y)$

\noindent $\varphi'$ = aplicația indusă pe $T$ de $\varphi$.

\noindent Fie $S = M_2(\mathbb{R})$, $(M_2(\mathbb{R}),\cdot)$

\[
T = \left\{ \begin{pmatrix} 1 & a \\ 0 & 1 \end{pmatrix} \,\middle|\, a\in\mathbb{R} \right\} \subseteq S = M_2(\mathbb{R}).
\]
\[
X_a X_b = X_{a+b}: \qquad \begin{pmatrix}1 & a\\ 0 & 1\end{pmatrix}\begin{pmatrix}1 & b\\ 0 & 1\end{pmatrix} = \begin{pmatrix}1 & a+b\\ 0 & 1\end{pmatrix}
\]

\bigskip
\noindent\fbox{Propoziție} \ Fie $S\neq\varnothing$, $\varphi$ lege de compoziție pe $S$.\\
$T\subseteq S$, $T$ stabilă în raport cu $\varphi$,\\
$\varphi'$ = operația indusă de $\varphi$ pe $T$: $\varphi'(x,y)=\varphi(x,y)\in T$.

\begin{enumerate}
\item[(1)] Dacă $\varphi$ este asociativă (resp.\ comutativă) $\Rightarrow$ $\varphi'$ este asociativă (resp.\ comutativă).
\item[(2)] Dacă $\varphi$ admite element neutru $e$ şi $e\in T$ $\Rightarrow$ $e$ este element neutru şi pentru $\varphi'$.
\item[(3)] Dacă un element $x\in T$ este simetrizabil în raport cu $\varphi$ ($\varphi$ = asociativă şi cu element neutru) şi simetricul său $x'$ în raport cu $\varphi$ este în $T$ $\Rightarrow$ $x$ este simetrizabil şi în raport cu $\varphi'$.
\item[(4)] simetricul (lui $x$ în raport cu $\varphi'$), fiind tot $x'$.
\end{enumerate}

\end{document}
